% Supplementary Information of the ERL Letter (D-107): robustness checks (brief section 10;
% verdict rule D-092; runs D-097). Every number is a \metric{} macro.

Each check was run as a variant of the primary analysis and judged by a fixed rule. A variant is ``robust'' if no headline decision changes and no magnitude moves materially (the elevation model's NMAD or bias, on the same points, beyond the tolerances fixed in the rule); ``partly sensitive'' if a sector-level outcome or a magnitude changes but no headline decision; and ``sensitive'' if a headline decision of H1 or H2 changes, in either direction. A check takes the worst verdict of its variants. The rule was finalised after the tables of two variants (fixed MNDWI~0 and AWEI\textsubscript{sh}) had been seen; the magnitude criterion added then can only make verdicts stricter.

Several variants (checks 1, 2, 3, 7 and 9) were run only on a seeded sample of the \metric{params.block_size_km}~km test blocks (sample R), and they were judged on the waterline elevation model: on the same R points, with the primary datum constant and nothing refitted, the headline is whether the model's NMAD stays below every grid's. These variants test the observation pipeline through the elevation model, not through the hindcast audit; the audit was recomputed only for the variants of checks 6 and 10. No check covers the near-term analysis (H3), whose sensitivity is the red-team test R3 (section~\SIredteam).

\begin{table}[!htbp]
\centering
\footnotesize
\setlength{\tabcolsep}{3pt}
\caption{\textbf{Robustness checks.} Verdict of each check (the worst of its variants) and the evidence; R, the seeded sample of test blocks; NMAD and bias of the waterline elevation model against ICESat-2 ATL06 unless stated. Variants not run are listed in the text.}
\label{tab:si_robust}
\begin{tabular}{>{\raggedright\arraybackslash}p{0.4cm}>{\raggedright\arraybackslash}p{3.6cm}>{\raggedright\arraybackslash}p{1.9cm}>{\raggedright\arraybackslash}p{9.0cm}}
\hline
\# & Check & Verdict & Evidence \\
\hline
1 & Per-acquisition threshold: Otsu against fixed MNDWI~0 and AWEI\textsubscript{sh} (R) & \metric{robustness.check01.verdict} & Bias on R moves from \metric{robustness.check01.fixed_mndwi0_R.atl06.dem_bias_m_primary}~m to \metric{robustness.check01.fixed_mndwi0_R.atl06.dem_bias_m}~m (MNDWI~0) and \metric{robustness.check01.awei_sh_R.atl06.dem_bias_m}~m (AWEI\textsubscript{sh}); NMAD (\metric{robustness.check01.awei_sh_R.atl06.dem_nmad_m}~m with AWEI\textsubscript{sh}) stays below every grid's (at least \metric{robustness.check01.awei_sh_R.atl06.min_product_nmad_m}~m) \\
2 & Sensor: without TM, ETM+, SLC-off or OLI scenes; Sentinel-2 added (R) & \metric{robustness.check02.verdict} & Without OLI the NMAD rises from \metric{erl.text.m3.robustness.check02.sub_oli_R.atl06.dem_nmad_m_primary} to \metric{robustness.check02.sub_oli_R.atl06.dem_nmad_m}~m and the model loses its lead on the northeastern coast; without TM ``\metric{robustness.check02.sub_tm_R.verdict}'', Sentinel-2 added ``\metric{robustness.check02.s2_add_R.verdict}''; the model leads overall in every variant \\
3 & Season and ice: the snow test in every month (R) & \metric{robustness.check03.verdict} & NMAD \metric{erl.text.m3.robustness.check03.snow_all_months_R.atl06.dem_nmad_m_primary} $\rightarrow$ \metric{robustness.check03.snow_all_months_R.atl06.dem_nmad_m}~m \\
4 & Wind setup: a fixed stricter threshold (R); setup correction instead of screening; screening in Komsomolets Bay & \metric{robustness.check04.verdict} & Stricter threshold ``\metric{robustness.check04.tau015_R.verdict}''; setup correction ``\metric{robustness.check04.setup_correction.verdict}'' (ATL08 NMAD \metric{robustness.check04.setup_correction.atl08.dem_nmad_m}~m); Komsomolets screening ``\metric{robustness.check04.komsomolets_screening.verdict}'': it lowers the waterline scatter at one threshold only and raises it at every stricter one \\
5 & Altimetry series and interpolation: Hydroweb with a constant offset; DAHITI by monotone cubic interpolation; a one-year window & \metric{robustness.check05.verdict} & ATL08 NMAD \metric{robustness.check05.hydroweb_const.atl08.dem_nmad_m}~m with Hydroweb levels, H2 unchanged; the DAHITI variants shift the scene levels by less than the rule's level tolerance \\
6 & Datums: a shape constant per sector (red-team test R2) & \metric{robustness.check06.verdict} & Sensitive in the strengthening direction: the SRTM15+ block test becomes significant ($p$ \metric{erl.text.p4.robustness.check06.sector_shape_constants.srtm15p_v2.7.block_test_p_primary_fix} $\rightarrow$ \metric{erl.text.p4.robustness.check06.sector_shape_constants.srtm15p_v2.7.block_test_p}) and every grid's majority outcome holds \\
7 & Interpolation: \metric{params.bin_width_sens_lo_m}~m bins (R); no static-contour exclusion & \metric{robustness.check07.verdict} & \metric{params.bin_width_sens_lo_m}~m bins raise the NMAD on R to \metric{robustness.check07.bins005_R.atl06.dem_nmad_m}~m; without the exclusion H2 stays not supported (ATL08 NMAD \metric{robustness.check07.no_static_exclusion.atl08.dem_nmad_m}~m) \\
8 & Validation block size: \metric{params.block_size_sens_lo_km} and \metric{params.block_size_km}~km & \metric{robustness.check08.verdict} & At \metric{params.block_size_sens_lo_km}~km H2 becomes supported for ATL08 (section~\SIdem); at \metric{params.block_size_km}~km the status holds but the supporting sectors change \\
9 & Vegetation: NDVI thresholds below and above the primary \metric{params.ndvi_veg} (R); reed beds as a population & \metric{robustness.check09.verdict} & The lower threshold raises the NMAD on R to \metric{robustness.check09.ndvi020_R.atl06.dem_nmad_m}~m and the model loses its northeastern lead; the higher one is ``\metric{robustness.check09.ndvi040_R.verdict}''; in reed beds the model is no better than the grids (section~\SIdem) \\
10 & Morphology and epochs: rising against falling limb; the audit on E2 and E3 only & \metric{robustness.check10.verdict} & The grids fit the older seabed better at \metric{erl.text.robust.morph_share_e1_better_pct}\% of shared levels (``\metric{robustness.check10.morphology_table.verdict}''); on E2 and E3 only every grid's majority outcome and block-test significance hold, and ETOPO 2022 rises to \metric{robustness.check10.h1_e2e3_only.etopo_2022.n_sectors_mispredict} sectors \\
11 & Prior work: the area--level slope of \citet{akbari2020} & \metric{robustness.check11.verdict} & Published \metric{lit.akbari2020.nc_area_per_m_km2}~\kmsq{} per metre; over the same levels the elevation model gives \metric{erl.text.k0.robustness.check11.akbari2020_area_slope.dem_band_range_-29.5_-26.5_km2_per_m}, \metric{erl.text.k0.robustness.check11.akbari2020_area_slope.dem_band_range_-29.0_-26.0_km2_per_m} and \metric{erl.text.k0.robustness.check11.akbari2020_area_slope.dem_band_range_-28.5_-25.5_km2_per_m}~\kmsq{} per metre under three datum readings of their levels: close under one, higher under the other two \\
\hline
\end{tabular}
\end{table}

\emph{Further variants.} Emergence dates shifted by one year: ``\metric{robustness.named.emergence_lag_365d.verdict}''; grids resampled by nearest neighbour (validation side): ``\metric{robustness.named.products_nearest.verdict}''.

\emph{Not run.} Gaussian-process (kriging) interpolation, \metric{params.bin_width_sens_hi_m}~m bins, the validation side of check~6 (constants at $\pm1\sigma$ and the mean-tide recipes), the as-dry and as-water bounds of check~9, the morphology map, the November season, the windows at the end of E1 and the reassignment of the calibration split were scripted but not run: memory and time ran out with three hindcasts in parallel. The audit sides of the nearest-neighbour and check~9 variants each need a full hindcast and were not run either.

\emph{Data quality.} The lon--lat edges of the export tiles bow in the equal-area grid, so \metric{qc.seam_gaps.unseen_km2}~\kmsq{} of the computation cover lies in neither neighbouring tile and is never observed; all of it lies within a few hundred metres of a tile seam. It is not fixed in this version.
